\section{Formulation and set-up}\label{sec:setup}

\subsection{Elements}\label{sec:elements}
The three elements use the same geometry, the same nodes and the same test functions:
Lagrange polynomials $N_i(\bxi)\in Q_p$ on the $(p+1)^2$ Gauss--Lobatto--Legendre
(GLL) nodes of the reference square, mapped isoparametrically, $\bx(\bxi)=\sum_i
N_i(\bxi)\,\bx_i$. They differ in the trial functions.
\begin{itemize}
\item \PM\ (unsymmetric spectral element of the code base, \texttt{Ke\_2D.m}):
  the trial functions are \emph{metric} functions $M_i^e(\bx)\in
  Q_p(\hat x,\hat y)$, defined element by element through $M_i^e(\bx_j)=\delta_{ij}$
  in the corotational skew coordinates $(\hat x,\hat y)^{\mathsf T}=\bm
  J_0^{-\mathsf T}(\bx-\bx_0)$ ($\bm J_0$, $\bx_0$: Jacobian and image of the element
  centre)~\cite{Xie2016Modified}. As $(\hat x,\hat y)$ is affine in $\bx$, the
  span of $\{M_i^e\}$ contains $P_p(\bx)$ (all polynomials of total degree $\le p$
  in Cartesian coordinates) for \emph{every} element shape. With
  $\bm B_L^\ast=\bm B_L\det\bm J$ (built from $\operatorname{adj}\bm J\,\nabla_{\bxi}N_i$),
  \begin{equation}
    \bm K^e=\int_{-1}^{1}\!\!\int_{-1}^{1}(\bm B_L^\ast)^{\mathsf T}\bm D\,\bm B_R\dd\xi\dd\eta ,
    \label{eq:Kpm}
  \end{equation}
  i.e.\ the Jacobian determinant cancels and the integrand is a polynomial in
  $\bxi$; this is the Jacobian-unlimited formulation PM-QUAD8$^\ast$ of
  \citet{Rajendran2010technique} generalised to degree $p$ (conditions C1--C3 of
  that paper). The trial space is non-conforming (the $M_i^e$ are not continuous
  across element edges). The code base solves for the $M_i^e$ in the Lagrange basis on the
  GLL points of the bounding box of the skew node coordinates
  (\texttt{p.metricBasis = 'lagrange'}, default); this changes the conditioning $\kappa(\bm V)$
  of the interpolation matrix, not the trial space, so all errors are those of the monomial
  basis up to round-off.
\item \PH\ (unsymmetric spectral element with hybrid trial space,
  \texttt{p.metricMethod = 2}): same test functions, Jacobian cancellation and
  \eqref{eq:Kpm}, but the trial functions span
  \begin{equation}
    T_p^e=P_p(\hat x,\hat y)\oplus\operatorname{span}\{L_a(\xi)L_b(\eta):\ a,b\le p,\ a+b>p\}
    \label{eq:hybrid}
  \end{equation}
  ($L_a$: Legendre polynomials), i.e.\ metric polynomials only for the complete space
  $P_p(\bx)$, which guarantees completeness and polynomial exactness, and parametric
  functions for the remaining $(p+1)^2-\tfrac12(p+1)(p+2)$ functions. $\dim T_p^e=(p+1)^2$; on
  bilinear elements $T_p^e=Q_p(\bxi)$, and on affine elements $T_p^e=Q_p(\bx)$
  (\texttt{docs}, chapter on the metric basis). The strains of the complement need
  $\bm J^{-1}$, so \PH\ is undefined where $\det\bm J=0$ at a quadrature point (not encountered
  here).
\item \PP\ (conventional spectral element, reference): parametric functions as
  test and trial functions, $\bm K^e=\int\bm B_L^{\mathsf T}\bm D\bm B_L\det\bm
  J\dd\xi\dd\eta$ (\texttt{tests/helpers/usemKeSymmetric.m}).
\end{itemize}
The load vectors are identical for the three elements (parametric test functions),
$\bm f^e=\int N_i\,\bm b\det\bm J\dd\xi\dd\eta+\int_{\Gamma_N}N_i\,\bar{\bm t}\dd s$.
The stiffness matrices use the default rule of \texttt{Ke\_2D.m}, $2p\times2p$
Gauss points; loads and error integrals use $(3p+2)^2$ points. The element
matrices of the study (\texttt{cyl\_Ke\_PM.m}, \texttt{cyl\_Ke\_PP.m}) are
vectorised versions of \texttt{Ke\_2D.m} and \texttt{usemKeSymmetric.m}; they agree
with them to round-off (\cref{tab:verification}). \texttt{ShapeFunctions2D\_Metric} stops with an
error when the interpolation matrix of an element is singular to working precision
($\mathrm{rcond}<\varepsilon$, both unsymmetric elements); such cases are marked ``rej.''.

\subsection{Model and mesh (gmsh)}\label{sec:mesh}
Thick cylinder under internal pressure (Section~5.9 of
\citep{Rajendran2010technique}): quarter annulus $a=1\le r\le b=3$, $E=1000$,
$\nu=0.25$, plane stress, internal pressure $p_i=1000$, symmetry conditions
$u_x=0$ on $x=0$ and $u_y=0$ on $y=0$. The geometry (exact circular arcs, built-in
kernel) and the structured mesh are generated with gmsh~\cite{Geuzaine2009Gmsh}
(\texttt{mesh/gen\_cylinder\_meshes.py}): a transfinite mesh of $2N\times2N$
eight-node quadrilaterals, i.e.\ $N\times N$ super-elements of $2\times2$ elements
($N=1,2,4,8$; element size $h\propto1/N$). Mesh distortion follows mesh
sequence~B of \citep[Fig.~19]{Rajendran2010technique}: in every super-element the
mid-side nodes of the four internal edges are rotated counter-clockwise about the
centre node by the angle $\alpha\in\{0^\circ,15^\circ,30^\circ,45^\circ,60^\circ,90^\circ\}$
(\cref{fig:meshes}). The pattern is repeated on every super-element, so the
\emph{relative} distortion does not decrease under refinement: the element maps are
not asymptotically affine. For $\alpha\ge60^\circ$ (and $\alpha=45^\circ$ for $N\ge2$),
$\det\bm J$ changes sign inside elements (folded elements, \cref{tab:folded}). Each mesh is written as
an ABAQUS input file (corner nodes first, node and element sets, pressure
\texttt{*Dsload}) that also runs with \texttt{code/Start\_Sim.m}.

The 8-node mesh is lifted to degree $p$ by \texttt{pSpectralElement\_NodesTopology.m}.
Element edges on $r=a$ and $r=b$ are then replaced by the exact arcs (Gordon--Hall
blending~\cite{Gordon1973Transfinite}); interior edges remain the parabolas of the
8-node mesh. The degree-$p$ geometry thus interpolates the circle at the GLL points;
the signed area converges to $2\pi$ spectrally ($|A_h-2\pi|<10^{-14}$ for $p=8$) and is
independent of $\alpha$ because folded parts cancel (\cref{tab:verification}).

\begin{figure}[t]
  \centering
  \begin{tikzpicture}
\begin{groupplot}[group style={group size=6 by 2, horizontal sep=1mm, vertical sep=3mm}, meshplot]
\nextgroupplot[title={$\alpha=0^\circ$}]
\addplot[black, line width=0.35pt] table {\pdat/mesh_N1_a000.dat};
\nextgroupplot[title={$\alpha=15^\circ$}]
\addplot[black, line width=0.35pt] table {\pdat/mesh_N1_a015.dat};
\nextgroupplot[title={$\alpha=30^\circ$}]
\addplot[black, line width=0.35pt] table {\pdat/mesh_N1_a030.dat};
\nextgroupplot[title={$\alpha=45^\circ$}]
\addplot[black, line width=0.35pt] table {\pdat/mesh_N1_a045.dat};
\nextgroupplot[title={$\alpha=60^\circ$}]
\addplot[black, line width=0.35pt] table {\pdat/mesh_N1_a060.dat};
\nextgroupplot[title={$\alpha=90^\circ$}]
\addplot[black, line width=0.35pt] table {\pdat/mesh_N1_a090.dat};
\nextgroupplot
\addplot[black, line width=0.35pt] table {\pdat/mesh_N2_a000.dat};
\nextgroupplot
\addplot[black, line width=0.35pt] table {\pdat/mesh_N2_a015.dat};
\nextgroupplot
\addplot[black, line width=0.35pt] table {\pdat/mesh_N2_a030.dat};
\nextgroupplot
\addplot[black, line width=0.35pt] table {\pdat/mesh_N2_a045.dat};
\nextgroupplot
\addplot[black, line width=0.35pt] table {\pdat/mesh_N2_a060.dat};
\nextgroupplot
\addplot[black, line width=0.35pt] table {\pdat/mesh_N2_a090.dat};
\end{groupplot}
\node[font=\scriptsize, rotate=90, anchor=south] at (group c1r1.west) {$N=1$};
\node[font=\scriptsize, rotate=90, anchor=south] at (group c1r2.west) {$N=2$};
\end{tikzpicture}

  \caption{Meshes generated with gmsh for $N=1$ ($2\times2$ elements) and $N=2$ ($4\times4$
    elements); the internal mid-side nodes of every super-element are rotated by $\alpha$.
    Meshes for $N=4,8$ repeat the pattern on $N^2$ super-elements.}
  \label{fig:meshes}
\end{figure}

\subsection{Problems}\label{sec:problems}
\begin{enumerate}
\item \emph{Lam\'e} (non-polynomial, rational): $u_r=\frac{a^2p_i\,r}{E(b^2-a^2)}\bigl[(1-\nu)+(1+\nu)\tfrac{b^2}{r^2}\bigr]$, $u_\theta=0$
  \cite[Eq.~(25)]{Rajendran2010technique}; symmetry conditions on the straight
  edges, tractions on the arcs.
\item \emph{Manufactured polynomial} field of total degree 5 ($X=x/3$, $Y=y/3$):\\
  $u_x=\tfrac1{10}(1+X+2Y+XY-3X^2Y+X^3+2X^4Y-X^2Y^3+X^5-2Y^5)$,\\
  $u_y=\tfrac1{10}(1-X+Y^2+2X^2Y-X^3+XY^3-3X^3Y^2+Y^5+2XY^4-X^5)$.
\item \emph{Manufactured smooth non-polynomial} (entire) field:
  $u_x=\tfrac1{10}\bigl(\sin3X\cos2Y+e^{XY}\bigr)$, $u_y=\tfrac1{10}\bigl(\cos2X\sin3Y+Ye^{-X}\bigr)$.
\end{enumerate}
For the manufactured solutions (MMS) both displacement components are prescribed on
$x=0$ and $y=0$ (nodal values) and the tractions $\bm\sigma(\bm u)\bm n_h$ on the
arcs. Strains, stresses and the body force $\bm b=-\operatorname{div}\bm\sigma(\bm u)$
are derived symbolically with the operators of the MMS module of the code base
(\texttt{MMS\_2D/differentialOperators}, as in \texttt{examples\_Elasticity.m});
for the Lam\'e field this yields $\bm b\equiv\bm0$. On the arcs $\bm n_h$ is the
normal of the discrete boundary, so that the discrete problem on $\Omega_h$ has
the exact field as solution whenever it lies in the trial space (with the exact
normal, the polynomial patch test fails at the level of the geometry error, e.g.\
$\eE\approx3\cdot10^{-6}$ for $p=5$, $N=1$).

\subsection{Error measures}\label{sec:errors}
With the strain field of the respective trial functions (metric for \PM, hybrid for \PH,
parametric for \PP),
\begin{equation}
  \eE^2=\frac{\sum_e\int_{\hat\Omega}(\bm\varepsilon-\bm\varepsilon_h^e)^{\mathsf T}\bm D
        (\bm\varepsilon-\bm\varepsilon_h^e)\det\bm J\dd\xi\dd\eta}
       {\int_\Omega\bm\varepsilon^{\mathsf T}\bm D\bm\varepsilon\dd\Omega},
  \qquad
  \eL^2=\frac{\sum_e\int_{\hat\Omega}|\bm u-\bm u_h^e|^2\det\bm J\dd\xi\dd\eta}{\int_\Omega|\bm u|^2\dd\Omega},
  \label{eq:errors}
\end{equation}
and the nodal error $e_n=\max_i|\bm u_h(\bx_i)-\bm u(\bx_i)|/\max_i|\bm u(\bx_i)|$.
The signed $\det\bm J$ is used, consistent with \eqref{eq:Kpm}: on folded elements
each element contributes with the degree of its map. For \PP\ on folded meshes the
parametric strain is singular where $\det\bm J=0$, so $\eE$ is not meaningful there;
$e_n$ is defined for all cases. Observed rates are $r=\log_2(e_N/e_{2N})$. For smooth
solutions and shape-regular, asymptotically affine meshes, optimal rates are $p$
($\eE$) and $p+1$ ($\eL$)~\cite{Ciarlet1972Interpolation}.
